\documentclass[11pt,paper=a4,answers]{exam}
\usepackage{graphicx,lastpage}
\usepackage{upgreek}
\usepackage{censor}
\usepackage{gensymb}
\usepackage{amsmath}
\usepackage{multicol}
\usepackage{enumitem}
\usepackage{tasks}
\usepackage{adjustbox}
\usepackage{xcolor}
\usepackage{tfrupee}
\setlength{\voffset}{0.25in}
\usepackage{tabularx}
\newcolumntype{Y}{>{\centering\arraybackslash}X}
%==============================================================
% Customise Non-Essential Packages Below
% =============================================================
\usepackage[most]{tcolorbox}
\usepackage{eso-pic}
\usepackage{tikz}
\AddToShipoutPictureBG{%
\begin{tikzpicture}[remember picture,overlay]
\foreach \x in {-8,-4,0,4,8,12,16,20}
\foreach \y in {-8,-4,0,4,8,12,16,20} {
\node[opacity=0.05,rotate=45,anchor=center] at ([xshift=\x cm,yshift=\y cm]current page.center) {%
\begin{minipage}{2cm}
\centering
\includegraphics[width=2cm]{images/logo_black.png}\\
\small preppeo.com
\end{minipage}%
};
}
\end{tikzpicture}%
}
\printanswers

%==============================================================
% Customise Non-Essential Packages Above
% =============================================================
\definecolor{svgray}{rgb}{0.90, 0.90, 0.90}
\graphicspath{{images/}}

\usepackage[normalem]{ulem}
\renewcommand\ULthickness{1pt}   %%---> For changing thickness of underline
\setlength\ULdepth{3.5ex}%\maxdimen ---> For changing depth of underline
\renewcommand{\baselinestretch}{1}
\chead{
\includegraphics[scale=0.1,height=2cm ]{logo.png}
}
\pagestyle{empty}

\pagestyle{headandfoot}
\headrule
\cfoot{W: https://preppeo.com; M: +91-8130711689}

\begin{document}
%==============================================================
% Customise Question paper details below this
% =============================================================
\begin{center}
    \centering
    \renewcommand{\arraystretch}{1.5}
    \begin{tabularx}{\textwidth}{YYY}
    & \normalsize\bf{{\Large{{Quant}}}} & \\
    By Preppeo & \normalsize\bf{{sat\_lid\_007}} & SAT\\
    \normalsize{{\textbf{{Unit:}} Algebra}} & \normalsize{{\textbf{{Chapter:}}Systems of two linear equations in two variables}} & \normalsize{{\textbf{{Topic:}} Solving Systems (Substitution \& Elimination)}} \\
    \end{tabularx}
\end{center}
\par\noindent\rule{\textwidth}{1pt}
% =============================================================
% Customise Question paper details above this
% =============================================================

\begin{questions}
\pointsinrightmargin
\pointsdroppedatright
\marksnotpoints
\pointpoints{mark}{marks}
\pointformat{\boldmath\themarginpoints}
\bracketedpoints

% =============================================================
% Question Begin
% =============================================================
%Q: 1
% Category: Practice | Difficulty: Medium | Question Type: Multiple Choice
\question
A local bakery sells boxes of donuts and muffins. A box of donuts costs \$12 and a box of muffins costs \$15. On Saturday, the bakery sold a total of 40 boxes for \$540. If $d$ represents the number of boxes of donuts and $m$ represents the number of boxes of muffins, which system of equations represents this situation?
\begin{tasks}(2)
    \task $d + m = 40$; $12d + 15m = 540$
    \task $d + m = 540$; $12d + 15m = 40$
    \task $12d + 15m = 40$; $d + m = 540$
    \task $d + m = 40$; $15d + 12m = 540$
\end{tasks}

\begin{solution}
\textbf{Conceptual Explanation:} \\
When modeling a real-world scenario with a system of equations, we typically look for two different types of totals. The first is usually a "quantity total" (how many items in total), and the second is a "value total" (how much money or weight in total).
\\[10pt]
\textbf{Calculation and Logic:} \\
The problem states that the total number of boxes is 40. Since $d$ is donuts and $m$ is muffins, our first equation is:
\[d + m = 40\]
The second piece of information is the total revenue of \$540. We multiply the cost of each item by its quantity to find the total value:
\[12d + 15m = 540\]
This gives us the system in option A.

\vspace{15pt}
Answer: \colorbox{yellow!100}{A}
\end{solution}



%Q: 2
% Category: Practice | Difficulty: Medium | Question Type: Numeric Entry
\question
Consider the system of equations below:
\begin{align*}
x + y &= 15 \\
2x - y &= 6
\end{align*}
What is the value of $x$ in the solution to the system?

\begin{solution}
\textbf{Conceptual Explanation:} \\
The elimination method is most efficient here because the $y$ terms have coefficients that are opposites (+1 and -1). When we add the two equations together, the $y$ variable will be eliminated, allowing us to solve for $x$ directly.
\\[10pt]
\textbf{Calculation and Logic:} \\
Add the two equations:
\[(x + y) + (2x - y) = 15 + 6\]
\[3x = 21\]
Divide by 3:
\[x = 7\]
The $x$-coordinate of the solution is 7.

\vspace{15pt}
Answer: \colorbox{yellow!100}{7}
\end{solution}

%Q: 3
% Category: Practice | Difficulty: Hard | Question Type: Multiple Choice
\question
A nutritionist is mixing two types of juices. Juice A contains 20\% sugar and Juice B contains 5\% sugar. The final mixture should be 10 liters and contain 1.25 liters of sugar. If $a$ is the liters of Juice A and $b$ is the liters of Juice B, which system of equations models this?
\begin{tasks}(2)
    \task $a + b = 10$; $0.20a + 0.05b = 1.25$
    \task $a + b = 1.25$; $0.20a + 0.05b = 10$
    \task $a + b = 10$; $20a + 5b = 1.25$
    \task $0.20a + 0.05b = 10$; $a + b = 10$
\end{tasks}

\begin{solution}
\textbf{Conceptual Explanation:} \\
In mixture problems, we create one equation for the total volume and a second equation for the specific component (in this case, sugar) contained within that volume.
\\[10pt]
\textbf{Calculation and Logic:} \\
The total volume of juice must be 10 liters:
\[a + b = 10\]
The amount of sugar from Juice A is $20\%$ of $a$ ($0.20a$) and from Juice B is $5\%$ of $b$ ($0.05b$). The total sugar is 1.25 liters:
\[0.20a + 0.05b = 1.25\]
This matches the system provided in option A.

\vspace{15pt}
Answer: \colorbox{yellow!100}{A}
\end{solution}

%Q: 4
% Category: Practice | Difficulty: Hard | Question Type: Numeric Entry
\question
If $(x, y)$ is the solution to the system of equations below, what is the value of $y$?
\begin{align*}
y &= 3x - 4 \\
2x + y &= 16
\end{align*}

\begin{solution}
\textbf{Conceptual Explanation:} \\
Since the first equation is already solved for $y$, the substitution method is the best approach. We can replace the $y$ in the second equation with the expression $3x - 4$ from the first equation.
\\[10pt]
\textbf{Calculation and Logic:} \\
Substitute $3x - 4$ into the second equation:
\[2x + (3x - 4) = 16\]
Combine like terms:
\[5x - 4 = 16\]
Add 4 to both sides:
\[5x = 20 \Rightarrow x = 4\]
Now, find $y$ by plugging $x = 4$ back into the first equation:
\[y = 3(4) - 4 = 12 - 4 = 8\]

\vspace{15pt}
Answer: \colorbox{yellow!100}{8}
\end{solution}

%Q: 5
% Category: Practice | Difficulty: Medium | Question Type: Multiple Choice
\question
A landscaping company has two types of trucks. Type X carries 3 tons of soil and Type Y carries 5 tons. The company used a total of 12 trucks to deliver 44 tons of soil. Which system represents this?
\begin{tasks}(2)
    \task $x + y = 12$; $3x + 5y = 44$
    \task $x + y = 44$; $3x + 5y = 12$
    \task $3x + 5y = 12$; $x + y = 12$
    \task $x + y = 12$; $5x + 3y = 44$
\end{tasks}

\begin{solution}
\textbf{Conceptual Explanation:} \\
We define two variables for the number of trucks and set up equations based on the count of vehicles and their total carrying capacity.
\\[10pt]
\textbf{Calculation and Logic:} \\
Total number of trucks:
\[x + y = 12\]
Total tons of soil (capacity per truck $\times$ number of trucks):
\[3x + 5y = 44\]
Option A correctly represents both constraints.

\vspace{15pt}
Answer: \colorbox{yellow!100}{A}
\end{solution}

%Q: 6
% Category: Practice | Difficulty: Medium | Question Type: Numeric Entry
\question
In the system below, what is the value of $x + y$?
\begin{align*}
2x + 3y &= 12 \\
x - y &= 1
\end{align*}

\begin{solution}
\textbf{Conceptual Explanation:} \\
We first solve the system to find the individual values of $x$ and $y$, then add them together as requested by the question.
\\[10pt]
\textbf{Calculation and Logic:} \\
Multiply the second equation by 3 to use elimination:
\[3(x - y) = 3(1) \Rightarrow 3x - 3y = 3\]
Add this to the first equation:
\[(2x + 3y) + (3x - 3y) = 12 + 3\]
\[5x = 15 \Rightarrow x = 3\]
Plug $x = 3$ into $x - y = 1$:
\[3 - y = 1 \Rightarrow y = 2\]
Final value: $x + y = 3 + 2 = 5$.

\vspace{15pt}
Answer: \colorbox{yellow!100}{5}
\end{solution}



%Q: 7
% Category: Practice | Difficulty: Hard | Question Type: Multiple Choice
\question
A hotel has 60 rooms, some of which are singles and some are doubles. A single room costs \$90 per night and a double room costs \$140 per night. If the hotel is full and the total revenue is \$7,400, how many double rooms ($d$) does the hotel have?
\begin{tasks}(4)
    \task 20
    \task 40
    \task 30
    \task 25
\end{tasks}

\begin{solution}
\textbf{Conceptual Explanation:} \\
We set up a system where $s$ is the number of single rooms and $d$ is the number of double rooms.
\\[10pt]
\textbf{Calculation and Logic:} \\
Equation 1 (Total rooms): $s + d = 60 \Rightarrow s = 60 - d$ \\
Equation 2 (Total revenue): $90s + 140d = 7,400$ \\
Substitute $s$ into Equation 2:
\[90(60 - d) + 140d = 7,400\]
\[5,400 - 90d + 140d = 7,400\]
\[50d = 2,000 \Rightarrow d = 40\]

\vspace{15pt}
Answer: \colorbox{yellow!100}{B}
\end{solution}

%Q: 8
% Category: Practice | Difficulty: Medium | Question Type: Numeric Entry
\question
If $2x + y = 10$ and $x + 2y = 8$, what is the value of $x - y$?

\begin{solution}
\textbf{Conceptual Explanation:} \\
Instead of solving for $x$ and $y$ separately, we can look for a way to manipulate the equations to get the expression $x - y$ directly.
\\[10pt]
\textbf{Calculation and Logic:} \\
Subtract the second equation from the first:
\[(2x + y) - (x + 2y) = 10 - 8\]
\[2x + y - x - 2y = 2\]
\[x - y = 2\]
This shortcut allows us to find the answer without full elimination.

\vspace{15pt}
Answer: \colorbox{yellow!100}{2}
\end{solution}

%Q: 9
% Category: Practice | Difficulty: Hard | Question Type: Multiple Choice
\question
A merchant buys two types of coffee beans. Type A costs \$4 per pound and Type B costs \$7 per pound. The merchant buys a total of 30 pounds for \$165. How many pounds of Type B coffee did he buy?
\begin{tasks}(4)
    \task 12
    \task 15
    \task 18
    \task 20
\end{tasks}

\begin{solution}
\textbf{Conceptual Explanation:} \\
Let $a$ be pounds of Type A and $b$ be pounds of Type B.
\\[10pt]
\textbf{Calculation and Logic:} \\
$a + b = 30 \Rightarrow a = 30 - b$ \\
$4a + 7b = 165$ \\
Substitute:
\[4(30 - b) + 7b = 165\]
\[120 - 4b + 7b = 165\]
\[3b = 45 \Rightarrow b = 15\]

\vspace{15pt}
Answer: \colorbox{yellow!100}{B}
\end{solution}

%Q: 10
% Category: Practice | Difficulty: Medium | Question Type: Numeric Entry
\question
Find the value of $x$ for the system:
\begin{align*}
3x + 2y &= 18 \\
x - 2y &= 2
\end{align*}

\begin{solution}
\textbf{Conceptual Explanation:} \\
The $y$ coefficients are $2$ and $-2$. Adding these equations will eliminate $y$.
\\[10pt]
\textbf{Calculation and Logic:} \\
\[(3x + 2y) + (x - 2y) = 18 + 2\]
\[4x = 20 \Rightarrow x = 5\]

\vspace{15pt}
Answer: \colorbox{yellow!100}{5}
\end{solution}

%Q: 11
% Category: Practice | Difficulty: Medium | Question Type: Multiple Choice
\question
An exam has 50 questions. Students earn 4 points for each correct answer ($c$) and lose 1 point for each incorrect answer ($w$). If a student gets a score of 150 by answering all questions, which system represents this?
\begin{tasks}(2)
    \task $c + w = 50$; $4c - w = 150$
    \task $c + w = 150$; $4c - w = 50$
    \task $4c + w = 50$; $c - w = 150$
    \task $c + w = 50$; $4c + w = 150$
\end{tasks}

\begin{solution}
\textbf{Conceptual Explanation:} \\
The total questions answered is the sum of correct and incorrect answers. The score is determined by the points gained minus the points lost.
\\[10pt]
\textbf{Calculation and Logic:} \\
Total questions: $c + w = 50$. \\
Total points: $4(c) - 1(w) = 150 \Rightarrow 4c - w = 150$.

\vspace{15pt}
Answer: \colorbox{yellow!100}{A}
\end{solution}

%Q: 12
% Category: Practice | Difficulty: Hard | Question Type: Numeric Entry
\question
If $5x + 2y = 20$ and $x + y = 7$, what is the value of $3x$?

\begin{solution}
\textbf{Conceptual Explanation:} \\
First, we solve for $x$ by eliminating $y$, then we multiply that result by 3.
\\[10pt]
\textbf{Calculation and Logic:} \\
Multiply $x + y = 7$ by 2:
\[2x + 2y = 14\]
Subtract from the first equation:
\[(5x + 2y) - (2x + 2y) = 20 - 14\]
\[3x = 6\]
The question asks for $3x$, which is 6.

\vspace{15pt}
Answer: \colorbox{yellow!100}{6}
\end{solution}

%Q: 13
% Category: Practice | Difficulty: Medium | Question Type: Multiple Choice
\question
Two numbers have a sum of 25 and a difference of 7. What is the larger number?
\begin{tasks}(4)
    \task 16
    \task 18
    \task 9
    \task 11
\end{tasks}

\begin{solution}
\textbf{Conceptual Explanation:} \\
Let the numbers be $x$ and $y$. We form two equations based on their sum and difference.
\\[10pt]
\textbf{Calculation and Logic:} \\
$x + y = 25$ \\
$x - y = 7$ \\
Add the equations: $2x = 32 \Rightarrow x = 16$. \\
Plug $x = 16$ into $x + y = 25$: $16 + y = 25 \Rightarrow y = 9$. \\
The larger number is 16.

\vspace{15pt}
Answer: \colorbox{yellow!100}{A}
\end{solution}

%Q: 14
% Category: Practice | Difficulty: Medium | Question Type: Numeric Entry
\question
In the system $x = 2y$ and $x + y = 12$, what is the value of $y$?

\begin{solution}
\textbf{Conceptual Explanation:} \\
Since $x$ is explicitly defined in terms of $y$, we substitute $2y$ for $x$ in the second equation.
\\[10pt]
\textbf{Calculation and Logic:} \\
\[(2y) + y = 12\]
\[3y = 12 \Rightarrow y = 4\]

\vspace{15pt}
Answer: \colorbox{yellow!100}{4}
\end{solution}

%Q: 15
% Category: Practice | Difficulty: Medium | Question Type: Multiple Choice
\question
A farmer has chickens ($c$) and cows ($w$). There are 20 animals in total and they have 56 legs. Which system represents this?
\begin{tasks}(2)
    \task $c + w = 20$; $2c + 4w = 56$
    \task $c + w = 56$; $2c + 4w = 20$
    \task $c + w = 20$; $4c + 2w = 56$
    \task $2c + 4w = 20$; $c + w = 56$
\end{tasks}

\begin{solution}
\textbf{Conceptual Explanation:} \\
Chickens have 2 legs and cows have 4 legs. The system consists of an animal count equation and a leg count equation.
\\[10pt]
\textbf{Calculation and Logic:} \\
Animal count: $c + w = 20$. \\
Leg count: $2c + 4w = 56$.

\vspace{15pt}
Answer: \colorbox{yellow!100}{A}
\end{solution}

%Q: 16
% Category: Practice | Difficulty: Hard | Question Type: Numeric Entry
\question
Solve for $x$:
\begin{align*}
4x - 3y &= 11 \\
2x + 3y &= 7
\end{align*}

\begin{solution}
\textbf{Conceptual Explanation:} \\
Adding these equations will eliminate the $y$ variable.
\\[10pt]
\textbf{Calculation and Logic:} \\
\[(4x - 3y) + (2x + 3y) = 11 + 7\]
\[6x = 18 \Rightarrow x = 3\]

\vspace{15pt}
Answer: \colorbox{yellow!100}{3}
\end{solution}

%Q: 17
% Category: Practice | Difficulty: Medium | Question Type: Multiple Choice
\question
A theater sells adult tickets for \$10 and child tickets for \$6. One evening, 100 tickets were sold for \$840. How many adult tickets ($a$) were sold?
\begin{tasks}(4)
    \task 40
    \task 60
    \task 50
    \task 70
\end{tasks}

\begin{solution}
\textbf{Conceptual Explanation:} \\
Let $a$ be adult and $c$ be child tickets.
\\[10pt]
\textbf{Calculation and Logic:} \\
$a + c = 100 \Rightarrow c = 100 - a$ \\
$10a + 6c = 840$ \\
\[10a + 6(100 - a) = 840\]
\[10a + 600 - 6a = 840\]
\[4a = 240 \Rightarrow a = 60\]

\vspace{15pt}
Answer: \colorbox{yellow!100}{B}
\end{solution}

%Q: 18
% Category: Practice | Difficulty: Hard | Question Type: Numeric Entry
\question
If $y = 2x$ and $3x + 2y = 21$, what is the value of $x$?

\begin{solution}
\textbf{Conceptual Explanation:} \\
Substitute $2x$ for $y$ in the second equation.
\\[10pt]
\textbf{Calculation and Logic:} \\
\[3x + 2(2x) = 21\]
\[3x + 4x = 21\]
\[7x = 21 \Rightarrow x = 3\]

\vspace{15pt}
Answer: \colorbox{yellow!100}{3}
\end{solution}

%Q: 19
% Category: Practice | Difficulty: Medium | Question Type: Multiple Choice
\question
$x + y = 10$ and $2x + 2y = 20$. How many solutions does this system have?
\begin{tasks}(4)
    \task Exactly one
    \task Exactly two
    \task Infinitely many
    \task None
\end{tasks}

\begin{solution}
\textbf{Conceptual Explanation:} \\
If one equation is a direct multiple of the other, they represent the same line.
\\[10pt]
\textbf{Calculation and Logic:} \\
The second equation $2x + 2y = 20$ can be divided by 2 to get $x + y = 10$, which is identical to the first equation. Since they are the same line, every point on the line is a solution.

\vspace{15pt}
Answer: \colorbox{yellow!100}{C}
\end{solution}

%Q: 20
% Category: Practice | Difficulty: Hard | Question Type: Numeric Entry
\question
If $ax + 3y = 12$ and $2x + y = 4$ has infinitely many solutions, what is the value of $a$?

\begin{solution}
\textbf{Conceptual Explanation:} \\
For infinitely many solutions, the equations must be equivalent.
\\[10pt]
\textbf{Calculation and Logic:} \\
The second equation has $y$, while the first has $3y$. Multiply the second equation by 3:
\[3(2x + y) = 3(4) \Rightarrow 6x + 3y = 12\]
Comparing $6x + 3y = 12$ to $ax + 3y = 12$, we see $a = 6$.

\vspace{15pt}
Answer: \colorbox{yellow!100}{6}
\end{solution}

%Q: 21
% Category: Practice | Difficulty: Medium | Question Type: Multiple Choice
\question
A car rental costs \$50 per day plus \$0.20 per mile ($m$). A truck rental costs \$30 per day plus \$0.30 per mile. At what distance in miles is the cost for one day the same for both?
\begin{tasks}(4)
    \task 100
    \task 200
    \task 150
    \task 250
\end{tasks}

\begin{solution}
\textbf{Conceptual Explanation:} \\
We set the total cost equations equal to each other to find the distance where the costs intersect.
\\[10pt]
\textbf{Calculation and Logic:} \\
Car: $50 + 0.20m$ \\
Truck: $30 + 0.30m$ \\
Set them equal:
\[50 + 0.20m = 30 + 0.30m\]
\[20 = 0.10m\]
\[m = 200\]

\vspace{15pt}
Answer: \colorbox{yellow!100}{B}
\end{solution}

%Q: 22
% Category: Practice | Difficulty: Hard | Question Type: Numeric Entry
\question
In the system $y = x + 4$ and $y = -2x + 10$, what is the $x$-coordinate of the solution?

\begin{solution}
\textbf{Conceptual Explanation:} \\
Since both equations equal $y$, we can set them equal to each other.
\\[10pt]
\textbf{Calculation and Logic:} \\
\[x + 4 = -2x + 10\]
\[3x = 6 \Rightarrow x = 2\]

\vspace{15pt}
Answer: \colorbox{yellow!100}{2}
\end{solution}

%Q: 23
% Category: Practice | Difficulty: Medium | Question Type: Multiple Choice
\question
$x + y = 5$ and $x + y = 8$. How many solutions?
\begin{tasks}(4)
    \task One
    \task None
    \task Infinitely many
    \task Two
\end{tasks}

\begin{solution}
\textbf{Conceptual Explanation:} \\
Two lines with the same slope but different intercepts are parallel and never intersect.
\\[10pt]
\textbf{Calculation and Logic:} \\
Both equations have a slope of -1 ($y = -x + 5$ and $y = -x + 8$). Since they are parallel and do not overlap, there are no solutions.

\vspace{15pt}
Answer: \colorbox{yellow!100}{B}
\end{solution}

%Q: 24
% Category: Practice | Difficulty: Hard | Question Type: Numeric Entry
\question
If $2x - 3y = 7$ and $x + y = 6$, find $y$.

\begin{solution}
\textbf{Conceptual Explanation:} \\
We eliminate $x$ by multiplying the second equation by 2.
\\[10pt]
\textbf{Calculation and Logic:} \\
Multiply $x + y = 6$ by 2: $2x + 2y = 12$. \\
Subtract the first from the second:
\[(2x + 2y) - (2x - 3y) = 12 - 7\]
\[5y = 5 \Rightarrow y = 1\]

\vspace{15pt}
Answer: \colorbox{yellow!100}{1}
\end{solution}

%Q: 25
% Category: Practice | Difficulty: Medium | Question Type: Multiple Choice
\question
A system has the same slope and different $y$-intercepts. The graph shows:
\begin{tasks}(2)
    \task Intersecting lines
    \task Parallel lines
    \task Perpendicular lines
    \task Overlapping lines
\end{tasks}

\begin{solution}
\textbf{Conceptual Explanation:} \\
Lines that change at the same rate (same slope) but start at different points (different intercepts) will always remain the same distance apart.
\\[10pt]
\textbf{Calculation and Logic:} \\
The definition of parallel lines in the $xy$-plane is having identical slopes and distinct $y$-intercepts.

\vspace{15pt}
Answer: \colorbox{yellow!100}{B}
\end{solution}

% --- PART 2: TOPIC KNOWLEDGE EXPANSION (26-50) ---

%Q: 26
% Category: Practice | Difficulty: Hard | Question Type: Numeric Entry
\question
If $3x - y = 10$ and $x + 2y = 8$, what is the value of $10y$?

\begin{solution}
\textbf{Conceptual Explanation:} \\
We solve for $y$ using elimination, then multiply the result by 10.
\\[10pt]
\textbf{Calculation and Logic:} \\
Multiply the first equation by 2: $6x - 2y = 20$. \\
Add to the second: $(6x - 2y) + (x + 2y) = 20 + 8 \Rightarrow 7x = 28 \Rightarrow x = 4$. \\
Substitute $x = 4$ into $x + 2y = 8$: $4 + 2y = 8 \Rightarrow 2y = 4 \Rightarrow y = 2$. \\
Final value: $10(2) = 20$.

\vspace{15pt}
Answer: \colorbox{yellow!100}{20}
\end{solution}



%Q: 27
% Category: Practice | Difficulty: Hard | Question Type: Multiple Choice
\question
For what value of $k$ will the system below have no solution?
\begin{align*}
kx - 3y &= 4 \\
4x - 6y &= 7
\end{align*}
\begin{tasks}(4)
    \task 2
    \task 4
    \task 8
    \task -2
\end{tasks}

\begin{solution}
\textbf{Conceptual Explanation:} \\
A system has no solution if the lines are parallel (same slope) but have different intercepts.
\\[10pt]
\textbf{Calculation and Logic:} \\
The $y$ coefficients are $-3$ and $-6$. Multiply the first equation by 2 to match the $y$ terms:
\[2kx - 6y = 8\]
For the lines to be parallel, the $x$ coefficients must match: $2k = 4 \Rightarrow k = 2$.
With $k=2$, the system is $4x - 6y = 8$ and $4x - 6y = 7$. Since $8 \neq 7$, there is no solution.

\vspace{15pt}
Answer: \colorbox{yellow!100}{A}
\end{solution}

%Q: 28
% Category: Practice | Difficulty: Hard | Question Type: Numeric Entry
\question
If $x = y + 3$ and $x^2 - y^2 = 21$, what is the value of $x + y$?

\begin{solution}
\textbf{Conceptual Explanation:} \\
We use the algebraic identity for the difference of squares: $x^2 - y^2 = (x - y)(x + y)$.
\\[10pt]
\textbf{Calculation and Logic:} \\
From the first equation, $x - y = 3$. \\
Substitute this into the identity:
\[(x - y)(x + y) = 21\]
\[(3)(x + y) = 21\]
Divide by 3:
\[x + y = 7\]

\vspace{15pt}
Answer: \colorbox{yellow!100}{7}
\end{solution}

%Q: 29
% Category: Practice | Difficulty: Hard | Question Type: Multiple Choice
\question
A system of two linear equations has no solutions. One equation is $y = 0.5x + 3$. Which could be the other?
\begin{tasks}(2)
    \task $x - 2y = 10$
    \task $2x - y = 6$
    \task $x + 2y = 3$
    \task $0.5x + y = 3$
\end{tasks}

\begin{solution}
\textbf{Conceptual Explanation:} \\
The other equation must have the same slope (0.5) but a different $y$-intercept.
\\[10pt]
\textbf{Calculation and Logic:} \\
Test Option A: $x - 2y = 10 \Rightarrow -2y = -x + 10 \Rightarrow y = 0.5x - 5$. \\
This line has the same slope (0.5) and a different intercept (-5), so it is parallel and has no solution.

\vspace{15pt}
Answer: \colorbox{yellow!100}{A}
\end{solution}

%Q: 30
% Category: Practice | Difficulty: Hard | Question Type: Numeric Entry
\question
If $x + y = 10$ and $x - y = 4$, what is the value of $x^2 - y^2$?

\begin{solution}
\textbf{Conceptual Explanation:} \\
We can use the difference of squares identity: $x^2 - y^2 = (x + y)(x - y)$.
\\[10pt]
\textbf{Calculation and Logic:} \\
We are given both parts of the identity: $x + y = 10$ and $x - y = 4$.
\[x^2 - y^2 = (10)(4) = 40\]

\vspace{15pt}
Answer: \colorbox{yellow!100}{40}
\end{solution}

%Q: 31
% Category: Practice | Difficulty: Medium | Question Type: Multiple Choice
\question
At a cafe, 2 coffees and 3 teas cost \$14.50. 3 coffees and 2 teas cost \$15.50. What is the cost of one coffee?
\begin{tasks}(4)
    \task \$3.50
    \task \$2.50
    \task \$4.00
    \task \$3.00
\end{tasks}

\begin{solution}
\textbf{Conceptual Explanation:} \\
Let $c$ be coffee and $t$ be tea. We use elimination to isolate $c$.
\\[10pt]
\textbf{Calculation and Logic:} \\
1) $2c + 3t = 14.50$ \\
2) $3c + 2t = 15.50$ \\
Multiply (1) by 2 and (2) by 3: \\
$4c + 6t = 29.00$ \\
$9c + 6t = 46.50$ \\
Subtract: $5c = 17.50 \Rightarrow c = 3.50$.

\vspace{15pt}
Answer: \colorbox{yellow!100}{A}
\end{solution}

%Q: 32
% Category: Practice | Difficulty: Medium | Question Type: Numeric Entry
\question
If $y = 3x - 1$ and $y = x + 7$, what is the value of $x$?

\begin{solution}
\textbf{Conceptual Explanation:} \\
Set the two expressions for $y$ equal to each other.
\\[10pt]
\textbf{Calculation and Logic:} \\
\[3x - 1 = x + 7\]
\[2x = 8 \Rightarrow x = 4\]

\vspace{15pt}
Answer: \colorbox{yellow!100}{4}
\end{solution}

%Q: 33
% Category: Practice | Difficulty: Hard | Question Type: Multiple Choice
\question
In the $xy$-plane, the lines $y = 2x + 1$ and $y = ax + b$ are perpendicular and intersect at $(0, 1)$. What is the value of $a$?
\begin{tasks}(4)
    \task -0.5
    \task 2
    \task -2
    \task 0.5
\end{tasks}

\begin{solution}
\textbf{Conceptual Explanation:} \\
Perpendicular lines have slopes that are negative reciprocals of each other.
\\[10pt]
\textbf{Calculation and Logic:} \\
The slope of the first line is 2. The negative reciprocal of 2 is $-1/2$ or $-0.5$. Thus, $a = -0.5$.

\vspace{15pt}
Answer: \colorbox{yellow!100}{A}
\end{solution}

%Q: 34
% Category: Practice | Difficulty: Medium | Question Type: Numeric Entry
\question
If $2x + 3y = 13$ and $y = 3$, what is the value of $x$?

\begin{solution}
\textbf{Conceptual Explanation:} \\
Substitute the given value of $y$ into the first equation.
\\[10pt]
\textbf{Calculation and Logic:} \\
\[2x + 3(3) = 13\]
\[2x + 9 = 13\]
\[2x = 4 \Rightarrow x = 2\]

\vspace{15pt}
Answer: \colorbox{yellow!100}{2}
\end{solution}

%Q: 35
% Category: Practice | Difficulty: Hard | Question Type: Multiple Choice
\question
How many solutions does the system below have?
\begin{align*}
x - 2y &= 5 \\
-3x + 6y &= -15
\end{align*}
\begin{tasks}(4)
    \task None
    \task One
    \task Infinitely many
    \task Two
\end{tasks}

\begin{solution}
\textbf{Conceptual Explanation:} \\
Check if the equations are multiples of each other.
\\[10pt]
\textbf{Calculation and Logic:} \\
Multiply the first equation by -3:
\[-3(x - 2y) = -3(5) \Rightarrow -3x + 6y = -15\]
Since this is exactly the second equation, they are the same line and have infinitely many solutions.

\vspace{15pt}
Answer: \colorbox{yellow!100}{C}
\end{solution}

%Q: 36
% Category: Practice | Difficulty: Hard | Question Type: Numeric Entry
\question
If $x + y = 20$ and $0.5x + 0.2y = 7$, what is $x$?

\begin{solution}
\textbf{Conceptual Explanation:} \\
We eliminate $y$ by multiplying the second equation by 5.
\\[10pt]
\textbf{Calculation and Logic:} \\
Multiply $0.5x + 0.2y = 7$ by 5:
\[2.5x + y = 35\]
Subtract $x + y = 20$ from it:
\[(2.5x + y) - (x + y) = 35 - 20\]
\[1.5x = 15 \Rightarrow x = 10\]

\vspace{15pt}
Answer: \colorbox{yellow!100}{10}
\end{solution}

%Q: 37
% Category: Practice | Difficulty: Medium | Question Type: Multiple Choice
\question
A boat travels 24 miles downstream in 2 hours and 24 miles upstream in 3 hours. What is the speed of the current?
\begin{tasks}(4)
    \task 2 mph
    \task 10 mph
    \task 12 mph
    \task 4 mph
\end{tasks}

\begin{solution}
\textbf{Conceptual Explanation:} \\
Let $b$ be the boat speed and $c$ be the current. Downstream speed is $b+c$, upstream is $b-c$.
\\[10pt]
\textbf{Calculation and Logic:} \\
$b + c = 24 / 2 = 12$ \\
$b - c = 24 / 3 = 8$ \\
Subtract the equations: $2c = 4 \Rightarrow c = 2$.

\vspace{15pt}
Answer: \colorbox{yellow!100}{A}
\end{solution}

%Q: 38
% Category: Practice | Difficulty: Medium | Question Type: Numeric Entry
\question
If $x = 5$ and $2x - y = 4$, what is $y$?

\begin{solution}
\textbf{Conceptual Explanation:} \\
Plug the known $x$ into the equation.
\\[10pt]
\textbf{Calculation and Logic:} \\
\[2(5) - y = 4\]
\[10 - y = 4 \Rightarrow y = 6\]

\vspace{15pt}
Answer: \colorbox{yellow!100}{6}
\end{solution}

%Q: 39
% Category: Practice | Difficulty: Hard | Question Type: Multiple Choice
\question
The lines $y = mx + 2$ and $y = 3x + 5$ are parallel. What is $m$?
\begin{tasks}(4)
    \task 3
    \task -3
    \task 1/3
    \task 2
\end{tasks}

\begin{solution}
\textbf{Conceptual Explanation:} \\
Parallel lines must have identical slopes.
\\[10pt]
\textbf{Calculation and Logic:} \\
The slope of the second line is 3. For the first line to be parallel, $m$ must also be 3.

\vspace{15pt}
Answer: \colorbox{yellow!100}{A}
\end{solution}

%Q: 40
% Category: Practice | Difficulty: Hard | Question Type: Numeric Entry
\question
Solve for $x$: $y = 2x - 5$ and $y = -x + 10$.

\begin{solution}
\textbf{Conceptual Explanation:} \\
Set the two expressions for $y$ equal.
\\[10pt]
\textbf{Calculation and Logic:} \\
\[2x - 5 = -x + 10\]
\[3x = 15 \Rightarrow x = 5\]

\vspace{15pt}
Answer: \colorbox{yellow!100}{5}
\end{solution}

%Q: 41
% Category: Practice | Difficulty: Hard | Question Type: Multiple Choice
\question
System $x + 2y = 10$ and $x + 2y = k$ has infinitely many solutions. $k = $?
\begin{tasks}(4)
    \task 10
    \task 5
    \task 20
    \task 0
\end{tasks}

\begin{solution}
\textbf{Conceptual Explanation:} \\
Equations must be identical for infinite solutions.
\\[10pt]
\textbf{Calculation and Logic:} \\
To match the first equation exactly, $k$ must be 10.

\vspace{15pt}
Answer: \colorbox{yellow!100}{A}
\end{solution}

%Q: 42
% Category: Practice | Difficulty: Medium | Question Type: Numeric Entry
\question
If $y = 4$ and $x + y = 10$, find $x$.

\begin{solution}
\textbf{Conceptual Explanation:} \\
Simple substitution.
\\[10pt]
\textbf{Calculation and Logic:} \\
\[x + 4 = 10 \Rightarrow x = 6\]

\vspace{15pt}
Answer: \colorbox{yellow!100}{6}
\end{solution}

%Q: 43
% Category: Practice | Difficulty: Hard | Question Type: Multiple Choice
\question
Perpendicular to $y = -0.25x$?
\begin{tasks}(4)
    \task $y = 4x$
    \task $y = -4x$
    \task $y = 0.25x$
    \task $y = -0.25x$
\end{tasks}

\begin{solution}
\textbf{Conceptual Explanation:} \\
Negative reciprocal of -0.25 (-1/4) is 4.
\\[10pt]
\textbf{Calculation and Logic:} \\
Slope of original line is $-1/4$. Flip and change sign: $4/1 = 4$.

\vspace{15pt}
Answer: \colorbox{yellow!100}{A}
\end{solution}

%Q: 44
% Category: Practice | Difficulty: Hard | Question Type: Numeric Entry
\question
If $3x + y = 13$ and $x + y = 7$, find $x$.

\begin{solution}
\textbf{Conceptual Explanation:} \\
Subtract the equations to eliminate $y$.
\\[10pt]
\textbf{Calculation and Logic:} \\
\[(3x + y) - (x + y) = 13 - 7\]
\[2x = 6 \Rightarrow x = 3\]

\vspace{15pt}
Answer: \colorbox{yellow!100}{3}
\end{solution}

%Q: 45
% Category: Practice | Difficulty: Medium | Question Type: Multiple Choice
\question
Sum is 10, difference is 2. Smaller number?
\begin{tasks}(4)
    \task 4
    \task 6
    \task 2
    \task 8
\end{tasks}

\begin{solution}
\textbf{Conceptual Explanation:} \\
$x+y=10, x-y=2$.
\\[10pt]
\textbf{Calculation and Logic:} \\
$2x=12 \Rightarrow x=6$. $6+y=10 \Rightarrow y=4$. Smaller is 4.

\vspace{15pt}
Answer: \colorbox{yellow!100}{A}
\end{solution}

%Q: 46
% Category: Practice | Difficulty: Hard | Question Type: Numeric Entry
\question
If $y = x^2$ and $y = 4$, how many positive solutions for $x$?

\begin{solution}
\textbf{Conceptual Explanation:} \\
Substitute $y$ into the first equation.
\\[10pt]
\textbf{Calculation and Logic:} \\
\[4 = x^2 \Rightarrow x = \pm 2\]
Positive solution is 2.

\vspace{15pt}
Answer: \colorbox{yellow!100}{1}
\end{solution}

%Q: 47
% Category: Practice | Difficulty: Medium | Question Type: Multiple Choice
\question
Intersection of $x=3$ and $y=5$?
\begin{tasks}(4)
    \task (3, 5)
    \task (5, 3)
    \task (0, 0)
    \task (3, 0)
\end{tasks}

\begin{solution}
\textbf{Conceptual Explanation:} \\
The intersection of a vertical line and a horizontal line is the point formed by their constant values.
\\[10pt]
\textbf{Calculation and Logic:} \\
$x$ is always 3, $y$ is always 5. Point is (3, 5).

\vspace{15pt}
Answer: \colorbox{yellow!100}{A}
\end{solution}

%Q: 48
% Category: Practice | Difficulty: Medium | Question Type: Numeric Entry
\question
If $y = 0$ and $2x + 3y = 12$, find $x$.

\begin{solution}
\textbf{Conceptual Explanation:} \\
Find the $x$-intercept.
\\[10pt]
\textbf{Calculation and Logic:} \\
\[2x + 3(0) = 12 \Rightarrow 2x = 12 \Rightarrow x = 6\]

\vspace{15pt}
Answer: \colorbox{yellow!100}{6}
\end{solution}

%Q: 49
% Category: Practice | Difficulty: Hard | Question Type: Multiple Choice
\question
Parallel to $y = 3x + 10$ through $(0,0)$?
\begin{tasks}(2)
    \task $y = 3x$
    \task $y = -3x$
    \task $y = 1/3x$
    \task $y = 3$
\end{tasks}

\begin{solution}
\textbf{Conceptual Explanation:} \\
Same slope, $y$-intercept is 0.
\\[10pt]
\textbf{Calculation and Logic:} \\
Slope must be 3. Through origin means $b=0$. $y = 3x$.

\vspace{15pt}
Answer: \colorbox{yellow!100}{A}
\end{solution}

%Q: 50
% Category: Practice | Difficulty: Hard | Question Type: Numeric Entry
\question
If $2x + y = 7$ and $x + 2y = 5$, find $x + y$.

\begin{solution}
\textbf{Conceptual Explanation:} \\
Add the two equations to find a multiple of $x + y$.
\\[10pt]
\textbf{Calculation and Logic:} \\
\[(2x + y) + (x + 2y) = 7 + 5\]
\[3x + 3y = 12\]
Divide by 3:
\[x + y = 4\]

\vspace{15pt}
Answer: \colorbox{yellow!100}{4}
\end{solution}


% =============================================================
% Question End
% =============================================================

\end{questions}
\begin{center}
\rule{.5\textwidth}{1pt}
\end{center}
\end{document}
